\section{FEMR binary result format}
\label{sec:femr-format}

FEMR version 1 is FEMaster's seekable binary result container. Multi-byte
integers and floating-point values are little-endian. Unknown chunks are
skipped using their stored size, allowing compatible format extensions.

Binary output is selected explicitly, for example:
\begin{verbatim}
FEMaster model.inp --output-format femr
FEMaster model.inp --output-format res frd femr
\end{verbatim}
FEMR version 1 always compresses each independently loadable field-data chunk
with LZ4. There is no compression setting in the command-line interface.

\subsection{Chunk envelope}

Every chunk begins with a 32-byte envelope containing, in order:
\begin{itemize}
  \item a four-byte chunk code;
  \item a \texttt{uint32} compression code (0 for none, 1 for an LZ4 block);
  \item \texttt{uint64} stored and uncompressed payload sizes;
  \item a \texttt{uint32} CRC32 of the uncompressed payload; and
  \item a reserved \texttt{uint32}, which is zero in version 1.
\end{itemize}
Only \texttt{FDAT} is compressed in version 1. Each field is a separate LZ4
block, so loading one field never decompresses another field.

\subsection{Chunks}

\begin{description}
  \item[HEAD] Must be first. It stores the \texttt{FEMR} magic, major and minor
  version, byte order, scalar width, default compression and reserved data.
  \item[INST] Stores each instance identifier once together with its
  length-prefixed UTF-8 name.
  \item[MESH] Starts with \texttt{uint64} node and element counts. Each node
  stores its dense/global, instance and part-local identifiers as three
  \texttt{int32} values, followed by three \texttt{float64} coordinates. Each
  element likewise stores its dense/global, instance and part-local identifiers,
  followed by a length-prefixed UTF-8 type name, node count and dense/global node
  identifiers.
  \item[LCAS] Stores the load-case identifier and step type.
  \item[FRAM] Stores load-case identifier, dense frame identifier and value.
  Eigenfrequency and buckling frame identifiers are derived from the numeric
  field-name suffix, so equal spectral values remain distinct frames.
  \item[FMET] Stores field identifier, load-case and frame identifiers, name,
  domain, scalar width, row count and component count.
  \item[FDAT] Stores the field identifier followed by its dense row-major scalar
  array.
  \item[CSUM] Must be last. Its payload is the CRC32 of every preceding chunk
  envelope and stored payload byte.
\end{description}

Strings use a \texttt{uint16} byte length followed by UTF-8 bytes. Version 1
supports only the field domains \texttt{UNKNOWN} (0), \texttt{NODE} (1), and
\texttt{ELEMENT} (2). Writers skip fields from other domains with a warning;
readers reject unsupported domain codes found in a version 1 file.

\subsection{Stable mesh and field indexing}

MESH ordering defines the stable dense result indexing. Node entry $i$ maps to
row $i$ of every NODE field, and element entry $i$ maps to row $i$ of every
ELEMENT field. Node and element identifiers are nevertheless stored explicitly;
they must not be reconstructed from element connectivity. The accompanying
instance and part-local identifiers preserve semantic labels such as
\texttt{17} and \texttt{bolt.17}.

\subsection{Lazy access}

A reader scans only chunk envelopes and the small metadata chunks. It seeks
over MESH and FDAT payloads and retains their offsets. Mesh or field payloads
are read, decompressed and checked only when requested.

\begin{verbatim}
import femaster_api as femaster

with femaster.open_results("model.femr") as results:
    stress = results.loadcase(1).frame(20).field("STRESS")
    mesh = results.mesh
\end{verbatim}
